\section{Universal bounds for positive multilinear functions}
\label{sec:positive-universal}
Throughout this section the factors are the positive monomials, with equal
supports collected. Let $R_d$ be the supremum of $T/H$ over all dimensions,
positive multilinear polynomials of maximum degree at most $d$, finite
nonnegative boxes, and points with $H>0$. Let $C_n$ be the corresponding
supremum in $n$ variables, with no additional degree restriction. Fixed
coordinates and affine summands may be removed. Proposition~\ref{prop:box-transfer}
shows that both suprema can equivalently be taken on the unit cube: expansion
does not increase dimension or degree and can only increase the termwise gap.
All logarithms without a subscript are natural.

\begin{theorem}[Sharp leading growth]\label{thm:positive-growth}
As the indicated integer parameter tends to infinity,
\begin{equation}\label{eq:positive-growth}
 R_d\sim\frac{\log d}{\log\log d},\qquad
 C_n\sim\frac{\log n}{\log\log n}.
\end{equation}
For every $d\ge2$, put $\Lambda=\max\{16,1+\log d\}$ and $c_0=1-e^{-1}$.
The explicit finite bound
\begin{equation}\label{eq:original-harmonic-finite}
 T\le\frac{24\Lambda}{c_0\log\Lambda}\,H
\end{equation}
holds on every finite nonnegative box. On the unit cube the upper bounds
are witnessed by distributions defined from the complete marginal vector
and degree allowance, independently of the supports and coefficients.
\end{theorem}

We first exhibit the obstruction to a constant bound, then construct a
common distribution that matches its order. The lower family disproves
Conjecture~1 in the author manuscript of \citet[p.~22]{Luedtke2012}, which
asks for a uniform constant for positive multilinear functions on
nonnegative boxes. The signed bilinear examples of Section~\ref{subsec:schur}
concern a different coefficient class.

\subsection{A sparse multiscale obstruction}
\begin{theorem}[Exact dyadic hull gap]\label{thm:dyadic-exact}
For every integer $L\ge2$, there is a unit-coefficient polynomial in
$n_L=2^L+L$ variables, with $2^{L+1}-2$ monomials, maximum degree
$d_L=2^{L-1}+1$, and an interior point such that
\begin{equation}\label{eq:dyadic-exact}
 T=L,\qquad H=s+(L-s)2^{-s},
\end{equation}
where $s\in\{1,\ldots,L-1\}$ satisfies
\begin{equation}\label{eq:dyadic-cutoff}
 (L-s+1)2^{-(s+1)}\le1\le(L-s+2)2^{-s}.
\end{equation}
Either adjacent choice at a cutoff equality gives the same value.
The total number of variable occurrences is $L2^L+2^{L+1}-2$.
\end{theorem}
\begin{proof}
Put $m=2^L$. Use leaf variables $z_1,\ldots,z_m$ and anchors
$a_1,\ldots,a_L$. For each $j$, partition the leaves into $2^j$ blocks
of size $k_j=2^{L-j}$, and define
\begin{equation}\label{eq:dyadic-polynomial}
 f_L(a,z)=\sum_{j=1}^L\sum_{B\in\mathcal P_j}
                   a_j\prod_{i\in B}z_i,
 \qquad a_j=2^{-j},\quad z_i=1-1/m.
\end{equation}
The upper and lower envelopes of each monomial at this point are $2^{-j}$
and $\max(0,2^{-j}-k_j/m)=0$. Thus $T=\cav f_L=L$.
The support and occurrence counts follow by summing $2^j$ and
$2^j(k_j+1)$, respectively.

For an arbitrary binary law with these means, let $R$ count failed leaves
and $N_j$ count blocks hit by a failure. Then $\E R=1$ and
$N_j\le\min(2^j,R)$. The deficiency representation gives
\begin{equation}\label{eq:dyadic-surrogate}
 H=\max\E\sum_j A_jN_j
 \le\sup\E\sum_j A_j\min(2^j,R),\qquad \E A_j=2^{-j}.
\end{equation}
For a fixed law of $R$, an anchor of mass $u$ maximizes its expectation
against any increasing function of $R$ by selecting the upper $u$-tail,
splitting an atom if necessary. Indeed, moving selected mass from a smaller
value to a larger unselected value cannot reduce the objective. A common
nonincreasing quantile $r(t)$, $0<t<1$, therefore makes every anchor
$A_j=\ind\{t\le2^{-j}\}$ optimal simultaneously. Consequently the right
side of \eqref{eq:dyadic-surrogate} equals the supremum of
$\int_0^1 S_{l(t)}(r(t))\,dt$, where
\[
 S_l(r)=\sum_{j=1}^l\min(2^j,r),\qquad
 w_l=\Prob(l(t)=l)=
 \begin{cases}2^{-l-1},&0\le l<L,\\2^{-L},&l=L.\end{cases}
\]
Here $r$ is integer valued, bounded by $m$, and has integral one.

Define $B_q=(L-q+2)2^{-q}$. The sequence decreases strictly, with
$B_1\ge1\ge B_L$, so an $s$ as in \eqref{eq:dyadic-cutoff} exists.
For every $r\ge0$,
\begin{equation}\label{eq:dyadic-affine-certificate}
 S_l(r)\le sr+F_l(s),\qquad
 F_l(s)=\begin{cases}0,&l\le s,\\2^{l-s+1}-2,&l>s.\end{cases}
\end{equation}
For $l\le s$, bound each term by $r$. For $l>s$, cap the first $l-s$
terms at $2^j$ and the remaining $s$ terms at $r$. Integrating yields
\[
 H\le s+\sum_lw_lF_l(s)=s+(L-s)2^{-s}.
\]

To attain this bound, choose $l$ with probabilities $w_l$, let
$A_j=\ind\{j\le l\}$, and for $q\in\{s,s+1\}$ put
\[
 R_l^{(q)}=\begin{cases}0,&l<q,\\2^{l-q+1},&l\ge q.\end{cases}
\]
Each anchor has its required mean and $\E R^{(q)}=B_q$. Mix $q=s$ and
$q=s+1$ with respective probabilities
$(1-B_{s+1})/(B_s-B_{s+1})$ and $(B_s-1)/(B_s-B_{s+1})$.
This gives $\E R=1$. For $l>s$ the two failure counts are the endpoints
of $[2^{l-s},2^{l-s+1}]$, where \eqref{eq:dyadic-affine-certificate}
is an equality. For $l=s$ they are $0,2$, also equality points, and for
$l<s$ both are zero.

It remains to realize the hit counts, a step requiring nested partitions.
Label leaves by $L$-bit strings, with level-$j$ blocks defined by the first
$j$ bits. The first $R$ strings in bit-reversal order have exactly
$\min(2^j,R)$ different prefixes of length $j$: their prefixes are the
reversed residues $0,\ldots,R-1$ modulo $2^j$. Declare these leaves failed
and apply a uniform bitwise XOR shift to every label. The shift preserves
all hit counts and gives each leaf failure probability $R/m$, conditional
on $R$. Averaging gives failure mean $1/m$ for every leaf. The constructed
law therefore attains the bound in the original polynomial.
\end{proof}

The cutoff conditions imply $s=\log_2 L+O(1)$ and $(L-s)2^{-s}=O(1)$.
Thus the ratio in \eqref{eq:dyadic-exact} is asymptotic to
$L/\log_2 L$. The exact formula uses nested dyadic partitions.
For arbitrary equal-size partitions the same surrogate proves the weaker
but still divergent lower bound
\begin{equation}\label{eq:arbitrary-partitions}
 \frac{T}{H}\ge\frac{L}{5+\log_2(1+(L-1)\log2)};
\end{equation}
Appendix~\ref{app:dyadic-alternatives} gives its distinct integral proof
and the dense symmetric predecessor.

The construction may be made homogeneous of degree $d_L$: introduce
$d_L-2$ padding variables and multiply a degree-$k$ term by the first
$d_L-k$ of them. At padding means one both gaps agree exactly with those
above, since all representing laws fix the padding coordinates. Moving
them sufficiently close to one makes all means interior with ratio
arbitrarily close to \eqref{eq:dyadic-exact}. This adds $O(n_L)$ variables
and keeps $O(n_L)$ monomials, but the occurrence count after padding is
$O(n_L^2)$, not the original $O(n_L\log n_L)$.

\subsection{Harmonic failure distributions}
Call a coordinate low when $x_i\le1/2$ and high otherwise. Independence
handles the terms with zero or at least two low coordinates by
Lemma~\ref{lem:easy-terms}. For the remaining terms we require failures
at several probability scales simultaneously.

Fix a real $M\ge d-1$, and put
\begin{equation}\label{eq:harmonic-parameters}
 \Lambda=1+\log M,\qquad \eta=(d-1)/M\in(0,1].
\end{equation}
We use $\eta$ for this cutoff parameter, to distinguish it from graph
density and physical box aspect ratio. Both laws below draw one uniform
$U\in(0,1)$ and set every low coordinate to $\ind\{U\le x_i\}$.
In the \emph{threshold law}, a high coordinate with failure mean
$p=1-x_i$ fails exactly when $U\le p$. In the \emph{harmonic law}, a
coordinate with $p=0$ never fails; otherwise set
\begin{equation}\label{eq:harmonic-law}
 h_p=\min(Mp,1),\quad A_p=1+\log(h_p/p),\quad
 q_p(t)=\frac{\min(1,p/t)}{A_p}\ind\{t\le h_p\}.
\end{equation}
Conditional on $U=t$, high coordinates fail independently with
probabilities $q_p(t)$. Since $M\ge1$, $1\le A_p\le\Lambda$ and
$h_p\ge p$. In particular these probabilities are valid, and
\[
 \int_0^1q_p(t)\,dt
 =\frac{p+p\log(h_p/p)}{A_p}=p.
\]
Thus the two laws preserve the full marginal vector, including boundary
means, and require no information about a particular term.

\begin{lemma}[Complete harmonic guarantee]\label{lem:harmonic-curve}
Consider a term with one low anchor of mean $u$, and high failure means
$p_1,\ldots,p_r$, where $r\le d-1$. Put
$t_e=\min(u,\sum_jp_j)>0$ and $z=\min(\max_jp_j/t_e,1)$. The threshold
and harmonic deficiencies are at least $zt_e$ and $J_{\Lambda,\eta}(z)t_e$,
respectively, where
\begin{equation}\label{eq:harmonic-integral}
 F_{\Lambda,\eta}(v)=1-\exp\left(-\frac{1-\eta v}{\Lambda v}\right),
 \qquad J_{\Lambda,\eta}(z)=\int_z^1 F_{\Lambda,\eta}(v)\,dv,
 \quad F_{\Lambda,\eta}(0)=1.
\end{equation}
\end{lemma}
\begin{proof}
Threshold deficiency is $\min(u,\max_jp_j)\ge zt_e$. If
$\max_jp_j\ge t_e$, the harmonic assertion is nonnegativity. Otherwise
integrate over $\max_jp_j\le t\le t_e$. The anchor succeeds there and
all $p_j\le t$. An inactive coordinate has $p_j<t/M$, so the inactive
failure mass is at most $rt/M\le\eta t$. Active mass is consequently at
least $\sum_jp_j-\eta t\ge t_e-\eta t$. On active coordinates,
$q_{p_j}(t)\ge p_j/(\Lambda t)$. Conditional independence gives union
failure probability at least
$1-\exp(-(t_e-\eta t)/(\Lambda t))$. Substitution $v=t/t_e$ proves the
claim. Zero term gaps require no division and satisfy every asserted
nonnegative bound automatically.
\end{proof}

\begin{theorem}[Optimized finite harmonic bound]\label{thm:harmonic-fixed-point}
For the parameters in \eqref{eq:harmonic-parameters}, let $\zeta$ be the
unique solution of $J_{\Lambda,\eta}(\zeta)=\zeta$ in $(0,1)$. Then
\begin{equation}\label{eq:harmonic-fixed-point-bound}
 R_d\le\frac1\zeta+\frac2{c_0}.
\end{equation}
If $w>0$ solves $we^w=\Lambda e^{-\eta}$ and $w\ge1$, then
\begin{equation}\label{eq:lambert-bound}
 \zeta\ge\frac{w-1/(2w)}{\Lambda},\qquad
 R_d\le\frac{\Lambda}{w-1/(2w)}+\frac2{c_0}.
\end{equation}
The parameter $M$ can be optimized over $[d-1,\infty)$ in either bound.
For $N=1+\log(d-1)$, choosing $M=(d-1)N$ gives, as $d\to\infty$,
\begin{equation}\label{eq:second-order-upper}
 R_d\le\frac{N}{\log N-\log\log N+o(1)}+\frac2{c_0}.
\end{equation}
This is an upper certificate, with no matching second-order lower assertion.
\end{theorem}
\begin{proof}
$F$ is continuous and strictly decreasing, positive on $(0,1)$, and
nonnegative at one. Thus $J$ is strictly decreasing and strictly convex,
and $J(z)-z$ has opposite signs at zero and one. Its unique root is
$\zeta$. Set $a=F(\zeta)\in(0,1)$. The convex tangent inequality gives
\[
 J(z)+az\ge J(\zeta)+a\zeta=(1+a)\zeta.
\]
Mix the harmonic and threshold laws with weights $1/(1+a)$ and
$a/(1+a)$. Lemma~\ref{lem:harmonic-curve} gives deficiency at least
$\zeta t_e$ for every one-low term. This mixture is optimal for the two
\emph{scalar guarantees} $J(z)$ and $z$: at $z=\zeta$ every convex
combination equals $\zeta$. It does not prove optimality over all actual
couplings or imply simultaneous attainability of these lower curves by
an extremal polynomial. Mix this law and independence with weights
proportional to $1/\zeta$ and $2/c_0$. Each term receives at least
$t_e/(1/\zeta+2/c_0)$ from its appropriate component; the other
contributions are nonnegative. Proposition~\ref{prop:deficiency} and
then Proposition~\ref{prop:box-transfer} prove
\eqref{eq:harmonic-fixed-point-bound}.

For the finite certificate, use $x-x^2/2\le1-e^{-x}\le x$ for $x\ge0$
and, for $0<z\le1$,
\[
 \int_z^1(1/v-\eta)^2\,dv
 =1/z-1+2\eta\log z+\eta^2(1-z)\le1/z.
\]
The last inequality follows from $2\eta\log z\le0$ and
$-1+\eta^2(1-z)\le0$. For $0<A\le\Lambda$ this gives
\begin{equation}\label{eq:lambert-integral-certificate}
 \Lambda J(A/\Lambda)
 \ge\log\Lambda-\log A-\eta+\eta A/\Lambda-1/(2A)
 \ge\log\Lambda-\log A-\eta-1/(2A).
\end{equation}
Take $A=w-1/(2w)$. Since $w\ge1$, $A\ge1/2$ and $A<\Lambda$.
Using $w+\log w=\log\Lambda-\eta$,
\[
 A+\log A+\frac1{2A}-(w+\log w)
 =\log\left(1-\frac1{2w^2}\right)+\frac1{4Aw^2}\le0.
\]
The inequality follows from $\log(1-t)\le-t$ and $A\ge1/2$.
Hence $J(A/\Lambda)\ge A/\Lambda$; monotonicity of $J(z)-z$ proves
\eqref{eq:lambert-bound}.

For the chosen cutoff, $\Lambda=N+\log N$ and $\eta=1/N$. The defining
identity for $w$ first gives
$\log\Lambda-\eta-\log\log\Lambda\le w\le\log\Lambda$
for all sufficiently large $N$. Thus $w/\log N\to1$, and substitution
back gives $w=\log N-\log\log N+o(1)$. Since
$(\log N)^2/N\to0$,
\[
 \frac{N}{\Lambda}\left(w-\frac1{2w}\right)
 =\log N-\log\log N+o(1).
\]
This proves \eqref{eq:second-order-upper}.
\end{proof}

\begin{proof}[Proof of Theorem~\ref{thm:positive-growth}]
For its original explicit finite bound, take
$\Lambda=\max\{16,1+\log d\}$ and $M=e^{\Lambda-1}\ge d$.
For a one-low term, put $t_e=\min(u,\sum_jp_j)>0$.
If $\max_jp_j\ge t_e/\sqrt\Lambda$, threshold deficiency is at least
$t_e/\sqrt\Lambda\ge t_e\log\Lambda/\Lambda$; here
$\log\Lambda\le\sqrt\Lambda$ follows by maximizing
$(\log t)/\sqrt t$. Otherwise integrate the harmonic law over
$[t_e/\sqrt\Lambda,t_e/2]$. Inactive failure mass is at most $t$,
so active mass is at least $t_e/2$. The conditional union probability is
at least $c_0t_e/(2\Lambda t)$, using $1-e^{-z}\ge c_0\min(1,z)$.
The harmonic deficiency is at least
\[
 \frac{c_0t_e}{4\Lambda}(\log\Lambda-2\log2)
 \ge\frac{c_0t_e\log\Lambda}{8\Lambda}.
\]
An equal mixture of independence, threshold and harmonic laws now gives
\eqref{eq:original-harmonic-finite}, since independence supplies at least
$c_0t_e/2$ on the other two classes. Every distribution is global.

Theorem~\ref{thm:harmonic-fixed-point} implies
$R_d\le(1+o(1))\log d/\log\log d$.
For the reverse inequality choose the largest $L$ with
$2^{L-1}+1\le d$ in Theorem~\ref{thm:dyadic-exact}. Its degree differs
from $d$ by a factor at most two, asymptotically, and its ratio is
asymptotic to $L/\log_2 L$. For dimension, $C_n\le R_n$; choose the
largest $L$ with $2^L+L\le n$ and add unused coordinates. Again the
dimensions differ by at most a constant factor. These comparisons prove
both equivalences in \eqref{eq:positive-growth} for all integers tending
to infinity, not merely for a dyadic subsequence.
\end{proof}

For comparison, the earlier leading-constant mixture has a distinct explicit
parameterization. Set $\Lambda=\max\{e^6,1+\log d\}$,
$M=e^{\Lambda-1}$, $b=\log\Lambda$, $a=\Lambda/b^2$, $\beta=1/b$, and
\begin{equation}\label{eq:older-leading-mixture}
 h=\frac{(1-1/b)(1-1/(2b^2))(b-3\log b)}{\Lambda}.
\end{equation}
Here $h>0$ because $b\ge6$ and $b-3\log b>0$.
If $\max p_j\ge t_e/a$, threshold deficiency is at least $t_e/a$.
Otherwise, on $[t_e/a,\beta t_e]$, active mass is at least
$(1-\beta)t_e$ and the lower exponent
$(1-\beta)t_e/(\Lambda t)$ is at most $1/b^2$.
Integrating $1-e^{-z}\ge z(1-1/(2b^2))$ gives harmonic deficiency
at least $ht_e$. Weights proportional to $1/h,a,2/c_0$ on harmonic,
threshold, and independence give $R_d\le1/h+a+2/c_0$, also with leading
constant one. The full-curve theorem retains more of the same distributions'
information and improves this certificate's lower-order loss.

A useful interpretation is simultaneous interpolation between Fr\'echet
bounds. The global law with $Z=1/\zeta+2/c_0$ satisfies, for every
$e\subseteq[n]$ with $2\le|e|\le d$,
\begin{equation}\label{eq:simultaneous-frechet}
 \E\prod_{i\in e}X_i\le(1-Z^{-1})\min_{i\in e}x_i
       +Z^{-1}\max(0,\sum_{i\in e}x_i-|e|+1).
\end{equation}
The best fraction guaranteed for every marginal vector is asymptotic to
$\log\log d/\log d$: the theorem gives the lower guarantee, and summing
any proposed guarantee over the dyadic example gives the reverse order
and constant. No convex-envelope evaluation algorithm is inferred.

\subsection{Removing coefficients without changing the degree supremum}
\begin{theorem}[Unit coefficients and interior homogeneous examples]
\label{thm:coefficient-removal}
For each fixed $d\ge2$, the supremum $R_d$ is unchanged by requiring
unit coefficients, degree exactly $d$ for every term, and strictly interior
unit-cube evaluation points simultaneously. Dimension and support size
are unrestricted in this assertion.
\end{theorem}
\begin{proof}
By positive expansion it suffices to approximate any unit-cube example.
Delete affine terms and homogenize with $d-2$ padding variables as above.
At padding means one the gaps are unchanged. Continuity of the finite
polyhedral envelopes (Lemma~\ref{lem:vertex-law}) permits an interior
perturbation preserving any strict target-ratio inequality. Normalize by
the largest coefficient. We now have one fixed homogeneous polynomial
$f=\sum_{e\in E}a_e\prod_{i\in e}x_i$, $0<a_e\le1$, with $n$ variables,
$s$ distinct supports, interior means $x$, and $H(f;x)>0$.

Replace coordinate $i$ by $m$ clones $y_{i1},\ldots,y_{im}$, and set
\[
 \bar y_i=m^{-1}\sum_{r=1}^my_{ir},\qquad F_m(y)=m^df(\bar y),
 \qquad y_{ir}=x_i.
\]
There are $sm^d$ distinct candidate clone monomials: each uses one clone
from each of its $d$ distinct original groups. At the repeated point,
\begin{equation}\label{eq:clone-envelopes}
 \vex F_m=m^d\vex f,\quad \cav F_m=m^d\cav f,\quad
 T(F_m)=m^dT(f),\quad H(F_m)=m^dH(f).
\end{equation}
Indeed, any binary clone law induces random group averages $q\in[0,1]^n$
with $\E q=x$. Conditional independent rounding of the original variables
with means $q_i$ has expected polynomial $f(q)$, so every normalized clone
expectation is an original feasible expectation. Conversely, setting all
clones in a group equal to its original binary coordinate embeds every
original law. This proves equality of the entire feasible expectation
intervals, and hence of both envelopes. Each clone term has its original
list of means, proving the termwise identity.

Independently retain each candidate term arising from $e$ with probability
$a_e$ and give it coefficient one; call the result $G_m$. For $N_0$
independent random variables with range lengths at most one, the bounded-sum
inequality of \citet[Theorem~2]{Hoeffding1963} gives
\begin{equation}\label{eq:bounded-sum}
 \Prob\bigl(\abs{\textstyle\sum_j(Y_j-\E Y_j)}>t\bigr)
 \le2e^{-2t^2/N_0}.
\end{equation}
For completeness, if $Y\in[a,b]$, the second derivative of
$\log\E e^{\lambda Y}$ is its variance under exponential tilting, at most
$(b-a)^2/4$: its squared distance from the interval midpoint is at most
that quantity. Integrating twice gives
$\E e^{\lambda(Y-\E Y)}\le e^{\lambda^2(b-a)^2/8}$.
Multiplication over independent variables, Markov's inequality, and
optimization at $\lambda=4t/N_0$ prove each tail in
\eqref{eq:bounded-sum}. Thus no concentration result is needed without
its hypotheses or proof.

Apply this inequality with $N_0=sm^d$ at each of the $2^{nm}$ clone
vertices, and once to the termwise gap at the repeated point. The latter
is a sum of the same retention indicators with fixed weights in $[0,1]$.
A union bound gives failure probability at most
\begin{equation}\label{eq:clone-union-bound}
 2(2^{nm}+1)\exp(-2t^2/(sm^d))
\end{equation}
for the two simultaneous assertions
\[
 \max_{v\in\{0,1\}^{nm}}|G_m(v)-F_m(v)|\le t,
 \qquad |T(G_m)-T(F_m)|\le t.
\]
Take $t=K m^{(d+1)/2}$ with $K^2>sn\log2/2$. The failure probability
tends to zero as $m\to\infty$, with $f,n,s,d$ fixed. Uniform vertex
error bounds the error of each envelope at every marginal vector by $t$,
by comparing expectations under all feasible laws. Thus the hull-gap
error is at most $2t$. Since $t/m^d\to0$ for $d\ge2$, successful samples'
normalized gaps converge to those of $f$, and their ratios converge
because its hull gap is positive. This proves equality of suprema, not
preservation at a fixed dimension or existence of a small sampled support.
\end{proof}
